\begin{theorem}[Wild odd residual coefficient table assembly]\label{mod:cases-wildodd-residualcoefficients}
Lean packages the completed lower and upper residual coefficient APIs as the
two table fields of a single public
\texttt{WildOddResidualCoefficientTable}.  It also exposes, by direct
reference to the already-proved theorems, the lower and upper
source-to-selected representative translations.  In particular, the polar
coefficient is unchanged, the affine translation has the positive
source-to-selected orientation, and the elementary stationary value is
transported together with the residual critical factor.

The same package keeps the two raw representative-ratio theorems separate
from the two indexed-character theorems.  Thus the low raw identity remains
\[
 \frac{r^{N,\mathrm{low}}_j}{r^{\mathrm{tw},\mathrm{low}}_j}
 =N_{K/F}\!\left(\frac{u+[j]}{[j]}\right)^{-1}z_j,
\]
and the corresponding high raw identity retains the inverse norm containing
the high source factor.  The manuscript $\mu_j(z_j)$ equality is exported
only after applying the indexed norm character and using its triviality on
that norm.

For the boundary, let $(\eta_0,\eta_0\gamma_0)$ be the named generator row,
let $(\eta,\eta\gamma)$ be the named base row, and put
\[
 d_0=\eta/\eta_0,
 \qquad \rho=d_0^{-1},\qquad
 \sigma=\gamma,\qquad \tau=\rho\gamma_0.
\]
Both named polar coefficients are nonzero, so Lean proves $d_0\ne0$.
The nonzero-index facts $[j]+d_0\ne0$ are exactly the boundary numerator
nonvanishing facts imported from the lower table; Lean adds the zero index
using $d_0\ne0$.  Define
\[
 a_0=\operatorname{Frob}^{-1}(1-d_0^{p-1}),
 \qquad
 b_0=\operatorname{Frob}^{-1}(\gamma-\gamma_0).
\]
The public \texttt{WildOddBoundaryAffinePencilCertificate} has exactly the
four proof fields consumed by \texttt{FiniteField.affinePencil}:
\[
 \rho\ne0,\qquad 1+[j]\rho\ne0
 \quad\text{for every prime-field index }j,
\]
and the exact identities
\[
 a_0^p=1-(\rho^{p-1})^{-1},
 \qquad b_0^p=\sigma-\rho^{-1}\tau.
\]
The principal assembly theorem simply installs the completed lower and upper
APIs, their representative transports, the raw and character-level ratio
statements, and this boundary constructor.  It does not construct a common
coordinate, reprove either coefficient table, invoke higher-symmetric
vanishing, apply the affine-pencil theorem, choose a stationary
representative, or perform phase cancellation.
\LeanFile{LanglandsFirstMainLemma/Cases/WildOdd/ResidualCoefficients.lean}
\lean{LanglandsFirstMainLemma.WildOddBoundaryAffinePencilCertificate}
\lean{LanglandsFirstMainLemma.wildOddBoundaryDZero_ne_zero}
\lean{LanglandsFirstMainLemma.wildOddBoundaryPrimeFieldShift_ne_zero}
\lean{LanglandsFirstMainLemma.wildOdd_boundaryAffinePencilCertificate}
\lean{LanglandsFirstMainLemma.WildOddResidualCoefficientTable}
\lean{LanglandsFirstMainLemma.wildOdd_twistCoefficients}
\leanok
\SourceText{Proposition prop:odd-coefficient-table, complete assembly and boundary certificate.}
\uses{mod:cases-wildodd-lowerresidualcoefficients,mod:cases-wildodd-upperresidualcoefficients}
\FormalizationContract
This assembly occupies fewer than 300 Lean lines and exports only the
completed tables, their translation and representative-ratio data, and the
four stated boundary facts.  The boundary proof is limited to prime-field
index packaging, nonzero division, and the defining inverse-Frobenius
equalities.  Any coefficient proof belongs in one of the two helper nodes,
and boundary phase cancellation belongs in BoundaryPencil.
\end{theorem}
